\section{Discussion}
\label{sec:discussion}

\noindent\textbf{What this enables for AI-assisted software engineering.}
The central positive result---that off-the-shelf LLMs can classify a kernel's \bandwidthbound/\computebound Roofline regime from software artifacts, without touching the target accelerator---is directly relevant to AI coding agents and developer-facing tools.
Before spending target-GPU time, an agent can ask a narrower question than ``will this be faster?'': is this kernel limited by data movement or by computation, and is the available evidence enough to answer that yet?
The answer is evidence-dependent, and numeric \rai serves as a secondary diagnostic: when it is accurate it reinforces the regime call, and when it falls in the heavy-tailed failure region it signals that the sample should be escalated to profiling.

\noindent\textbf{Where \sourceonly and \sourcesass fit in a real workflow.}
Execution-free reasoning has two meaningful stages.
\sourceonly fits pre-compilation tasks such as code review, design exploration, and agent planning; \sourcesass fits a later stage where the code can be compiled for a target architecture but the accelerator is still unavailable for execution or profiling.

\noindent\textbf{Why deterministic and learned static features still matter.}
The reference baselines sharpen the methodological boundary: a simple lexical rule cannot explain the useful FP32/FP64 \sourceonly signal, yet supervised source baselines recover corpus-level triage cues and architecture-specific SASS counts expose substantial post-compilation structure.
Our results suggest a cascade combining deterministic compiler or disassembly features, lightweight supervised models, and LLM reasoning could outperform any single strategy by routing each evidence tier to the method that wins it.
\autoref{tab:deployment-matrix} records the per-regime first-stage choice the released evidence supports---a routing guide, not a measured system, since we have not built the combined cascade and choosing a route presumes the precision and evidence tier are known in advance.

\begin{table}[t]
  \centering
  \caption{Deployment routing summary from the shared subset. These are corpus-backed first-stage choices, not universal rules.}
  \label{tab:deployment-matrix}
  \scriptsize
  \setlength{\tabcolsep}{3pt}
  \begin{tabular}{@{}p{0.18\columnwidth}p{0.22\columnwidth}p{0.46\columnwidth}@{}}
    Regime & First stage & Why \\
    \midrule
    \sourceonly FP16 & Static-first & One-shot LLMs stay at the trivial \bandwidthbound/\computebound baseline; cheap source features recover more queueing signal. \\
    \sourceonly FP32/FP64 & Frontier LLM & Best pre-compilation triage overall, especially for FP32 where paired LLM wins are decisive. \\
    \sourcesass FP16/FP32 & Frontier LLM & Largest paired gains once lowered arithmetic and DRAM-visible behavior become explicit. \\
    \sourcesass FP64 & Frontier LLM & Learned SASS features edge the label call, but not significantly, and their numeric \rai is several times coarser, so the frontier LLM stays the better single choice. \\
  \end{tabular}
\end{table}


\noindent\textbf{What tool builders should and should not take from the results.}
Tool builders should use the method as a triage filter---to prioritize profiling effort, annotate likely memory- versus compute-side cases, or flag kernels too uncertain to trust---not as a runtime substitute, speedup estimator, or universal static analysis of GPU performance.

\noindent\textbf{Why the benchmark release matters.}
Even with the method's limitations, the released corpus is a useful contribution in its own right: it exposes a benchmarkable software-engineering task with source files, build context, runtime arguments, SASS, profiler-derived labels, and model outputs, and the added deterministic and learned reference baselines make it a testbed for comparing execution-free reasoning strategies rather than only for ranking LLMs.
Triaging the entire 732-sample shared subset with the budget model costs under \$0.50, which keeps such comparisons cheap to run.
It supports several follow-on directions that require new experiments: stronger compiler-grade analyzers, repeated-query robustness studies, and hybrid predictors that combine compiler features with LLM reasoning.
